\section{Introduction}
\label{sec:intro}

% Problem
The wind that a truck meets on the road decides a large part of its energy use, yet it is almost never measured. At highway speeds aerodynamic drag is the largest single resistance term of a tractor--semitrailer, and it depends on the air velocity relative to the vehicle: head- and tailwind change the relative speed, and crosswind increases the yaw angle and, through the yaw dependence of the drag area, the drag coefficient \citep{Askerdal2024crosswind}. Over a return trip head- and tailwind partly cancel, crosswind does not. Range prediction for battery-electric trucks therefore needs the local wind along the route \citep{Askerdal2025winter,Shiledar2026}, and strong crosswind is in addition a safety issue for high-sided vehicles \citep{Baker1986,ArayaReyes2026}.

% Why the obvious data source fails
In practice the wind is taken from the nearest weather station or from a numerical weather prediction (NWP) product \citep{Filla2025data,Filla2026}. Both describe the wind at \SI{10}{m} above open, flat terrain, the standard exposure of synoptic observations \citep{WMO8}. A vehicle experiences the wind at \SIrange{2}{4}{m}, often in a corridor cut through forest, along a valley or between buildings. Vehicle measurements make the mismatch concrete. On Canadian highways bordered by forest or terrain, the on-road wind turned with the road and fell to about a fifth of the reference wind when the wind blew across the road \citep{McTavish2024}. In a winter test of a battery-electric truck, the cab-mounted anemometer measured \SI{3.7}{m.s^{-1}} where a station \SI{11}{km} away reported \SI{2.1}{m.s^{-1}} \citep{Askerdal2025winter}. The most complete long-term analysis of truck air resistance from recorded weather therefore lists the unknown height and exposure of the station wind as its main limitation \citep{Filla2026}.

% What others have done (recent)
Two research lines address parts of this problem. Meteorology and wind engineering correct station wind for exposure and roughness \citep{Wieringa1986,Wieringa1993,Verkaik2000}, and recent work has improved the roughness length itself, for example from anemometer towers \citep{Wang2024z0}; large-eddy simulations quantify how forest density, clearings and forest growth change the wind near forests \citep{Chaudhari2026,Topic2026}, and forest growth is a candidate cause of the decline of observed wind in Sweden \citep{Minola2023}. Statistical and deep-learning downscaling map coarse NWP wind to fine terrain \citep{HowardClark2007,Winstral2017,DujardinLehning2022,Dupuy2023,LeToumelin2023,LeToumelin2024,Lian2024}. These studies target wind at \SI{10}{m} or above, in mountains or over open land, evaluated at meteorological stations; a recent study reports that the speed correction learned in the Alps did not transfer to other mountain ranges \citep{LeToumelin2024}. On the vehicle side, on-road and on-bridge measurements characterise channelling and crosswind loads \citep{McTavish2024,Reymert2024,Reymert2025}, yaw-angle distributions have been linked to meteorological data \citep{Xiao2023}, and simplified drag models turn wind into force \citep{Askerdal2024crosswind}. None of these studies quantifies, out of sample and across regions, how station wind must be corrected to represent the wind at the side of a road, nor how much this matters for vehicle drag (Section~\ref{sec:related}).

% Opportunity
Road weather networks offer the missing intermediate reference. The Norwegian Public Roads Administration (Statens vegvesen) and the Swedish Transport Administration (Trafikverket) operate hundreds of roadside anemometers, exposed to the same corridors that vehicles drive through. They do not represent the wind at vehicle height along a whole road, but they are numerous, located on many different roads, and report continuously.

% This study
Here we ask what the wind costs a heavy truck on a given road, and how well this cost can be estimated from open data. We use \noTrv{} Norwegian roadside stations over two years, \noRef{} reference stations and open laser-scanning geodata, with a check in a second landscape in western Sweden, to answer three questions. How much, and why, does the wind change from a synoptic station to the roadside? Can a simple, physically interpretable correction remove most of this change for roads where it was not fitted? And what do the station wind and the corrected wind imply for the air-drag energy of a tractor--semitrailer? We show that stations overestimate the mean roadside wind by \noLadRawBiasAbs\,\%, that the difference is explained by the exposure of each station and by forest at two distinct scales, and that a correction combining all stations reduces the hourly wind error by \noCvRed\,\% out of sample. For the truck, the roadside wind adds on average \enAddNo\,\% to the air-drag energy, almost entirely through crosswind, but the penalty ranges from \enRoadMinNo{} to \enRoadMaxNo\,\% between roads. The nearest-station wind doubles the penalty and predicts the hourly energy worse than ignoring the wind, whereas the corrected wind reproduces the penalty of each road. Its main remaining weakness, the wind direction in mountain terrain, defines the next step.

% Outline
Section~\ref{sec:related} positions the study in recent literature, Section~\ref{sec:data} describes the regions and the data, Section~\ref{sec:methods} the correction chain and the validation design, Section~\ref{sec:results} the results, and Section~\ref{sec:discussion} discusses mechanisms, limits and the use for vehicles.
